% Generated draft: requires mathematical and source review.
\documentclass[11pt]{article}
\usepackage[margin=25mm]{geometry}
\usepackage{fontspec}
\setmainfont{Latin Modern Roman}
\usepackage{amsmath,amssymb,mathtools,unicode-math}
\setmathfont{Latin Modern Math}
\usepackage{enumitem,needspace,xurl,hyperref}
\setlength{\parindent}{0pt}
\setlength{\parskip}{5pt}
\setlength{\emergencystretch}{3em}
\newcommand{\N}{\mathbb{N}}
\newcommand{\Z}{\mathbb{Z}}
\newcommand{\Q}{\mathbb{Q}}
\newcommand{\R}{\mathbb{R}}
\newcommand{\C}{\mathbb{C}}
\newcommand{\F}{\mathbb{F}}
\newcommand{\A}{\mathbb{A}}
\newcommand{\PP}{\mathbb{P}}
\newcommand{\E}{\mathbb{E}}
\newcommand{\eps}{\varepsilon}
\newcommand{\dd}{\,\mathrm{d}}
\begin{document}
\section*{UnsolvedMath ID: 30002042}
\textbf{Source problem code:} catalogue entry
\subsection{Definitions and mathematical statement}
Let \(M\) be a lattice of rank \(d\), and let \(P\subset M\otimes_{\Z}\R\) be a full-dimensional lattice polytope, meaning the convex hull of finitely many lattice points with \(\dim P=d\). Lower-dimensional polytopes are understood in their affine lattice. Write \(M^{\vee}=\operatorname{Hom}(M,\Z)\). A lattice polytope \(Q\) containing the origin in its interior is reflexive if its polar
\[
Q^{\vee}=\{y\in M^{\vee}\otimes_{\Z}\R:\langle y,x\rangle\geq -1\text{ for every }x\in Q\}
\]
is a lattice polytope. A Gorenstein polytope of index \(r\in\Z_{>0}\) is a lattice polytope \(P\) such that \(rP\) has a unique interior lattice point \(m\) and \(rP-m\) is reflexive. Set
\[
n(P)=\dim(P)+1-2r.
\]
The supplied record denotes the stringy \(E\)-function of \(P\) by \(E_{\mathrm{st}}(P,t)\), but supplies neither its defining formula nor the convention relating this single-variable notation to any multivariable invariant. We retain this notation without selecting a specialization or normalization. For a nonzero polynomial \(f(t)=\sum_i a_i t^i\), its ordinary degree is \(\deg_t f=\max\{i:a_i\neq0\}\).

The source asks the following two questions for arbitrary Gorenstein polytopes:
\begin{enumerate}
\item For every Gorenstein polytope \(P\) of index \(r\), does its stringy \(E\)-function have the expected degree
\[
\deg_t E_{\mathrm{st}}(P,t)\stackrel{?}{=}\dim(P)+1-2r?
\]
\item For each fixed integer \(n\), are there only finitely many stringy \(E\)-polynomials in the family
\[
\bigl\{E_{\mathrm{st}}(P,t): P\text{ is a Gorenstein polytope of index }r,\ \dim(P)+1-2r=n\bigr\}
\]
up to multiples?
\end{enumerate}
Here dimension and index may vary while \(n\) remains fixed. The phrase ``up to multiples'' is retained from the source: the permitted multipliers and the resulting equivalence relation are unspecified. The record also gives no restriction to nonnegative \(n\), no treatment of an identically zero \(E\)-polynomial, and no degree convention for that case. Consequently, the displayed questions preserve the source formulation but require these qualifications to be resolved before becoming fully specified mathematical assertions.
\subsection{Short English statement}
The problem asks whether every Gorenstein polytope has a stringy \(E\)-polynomial of degree \(\dim(P)+1-2r\), and whether fixing this value leaves only finitely many such polynomials up to multiples. The invariant's variable convention, the meaning of multiples, and the treatment of zero polynomials and negative expected degrees remain unspecified. The supplied record labels the problem open; this draft does not independently verify that label.
\subsection{Sources}
\begin{itemize}
\item Supplied local registry: \texttt{data/unsolvedmath/problems.json}, permanent UnsolvedMath ID \(30002042\), source problem code \texttt{catalogue entry}, title ``Degrees and Finiteness of Stringy \(E\)-Polynomials.'' The record lists proposed year 2012 and names no proposer.
\item Supplied report citation: \emph{Toric Geometry (2013)}, Oberwolfach Reports, pp.\ 1231--1306, DOI \texttt{10.4171/owr/2012/21}. \url{https://doi.org/10.4171/owr/2012/21}. The supplied verification note states that the official report was checked and confirms the complete two-part question; its supplied statement status is \texttt{corrected\_verified}. These are source-record claims, not verification performed for this draft. The citation's 2013 label and the record's 2012 year are retained without reconciliation.
\item Nill--Schepers, \emph{Gorenstein polytopes and their stringy E-functions}. \url{https://doi.org/10.1007/s00208-012-0792-2}.
\item Nill--Schepers, \emph{On smooth Gorenstein polytopes}. \url{https://www.jstage.jst.go.jp/article/tmj/67/4/67_513/_pdf/-char/en}.
\item Literature assessment dated exactly \texttt{2026-08-21}: the supplied curation labels the problem open and reports that polynomiality is proved and finiteness is known for special smooth classes, while the expected-degree and general finiteness questions for arbitrary Gorenstein polytopes remain unresolved. This is supplied dated curation derived from the Oberwolfach Reports record, not independently verified literature assessment. The background explicitly requests independent verification of time-sensitive status and cited claims.
\end{itemize}
\end{document}
